% IK2215 Network Design Report Template
% Created by Voravit Tanyingyong <voravit@kth.se>
% Revised on 2023-08-25
%
% For how to use Overleaf and LaTeX, see Overleaf documentation below:
% https://www.overleaf.com/learn
%
% Page limit: maximum 6 pages!
% The TODO text contains a description of what you should consider in general and 
% some questions that you should ask yourself while designing the network. 
% They are there as a guideline, and you must remove them from your network design report!
% The report should be written in your own words describing your network design.
% Questions in TODO text must not be included as part of your report!
% 
\documentclass{article}
\usepackage[a4paper, total={6in, 9in}]{geometry}
\usepackage[T1]{fontenc}
\usepackage{tabto}
\usepackage{multirow}
\usepackage{array}
\usepackage{amsmath}
\usepackage{hyperref}
\usepackage[table]{xcolor}
\usepackage[color=yellow]{todonotes}
\usepackage{menukeys}
%\usepackage[utf8]{inputenc}
%\usepackage[english]{babel}
\usepackage{lastpage}
\usepackage{fancyhdr}
\usepackage{graphicx}
\usepackage{svg}
\pagestyle{fancy}
\fancyhf{}
\renewcommand{\headrulewidth}{0pt}
\cfoot{Page \thepage \hspace{1pt} of \pageref{LastPage}}
\newcommand\RED[1]{\textcolor{red}{#1}}

\title{IK2215: Network Design Report}
\author{
{Emanuel Paraschiv}\\
{emapar@kth.se}
\and
{Haoyu Chen}\\
{chchen3@kth.se}
}
% Just leave the date blank as it is. You do not need to have the date in the report.
\date{}


\begin{document}
\maketitle
\thispagestyle{fancy}

\section{General Information}
\textbf{ASN:102} 
\tab
\textbf{NETWORK:1.102.0.0/20} 

\section{Network overview}
This section contains an overview of network design.



\subsection{Network diagram}
The network design is illustrated in \figureautorefname~\ref{fig:your_figure} below.
\begin{figure}[h]
    \centering
    \includesvg[width=0.82\linewidth]{scheme}
    \caption{Network Diagram}
    \label{fig:your_figure}
\end{figure}

\subsection{IP address allocation}
The IP address allocation and domain name of each device in AS102 are presented in the table below:

\begin{center}
\begin{tabular}{ |l|l|l|l| } 
\hline
\textbf{Device} & \textbf{Interface} & \textbf{IP address} & \textbf{Domain name} \\
\hline
r1 & eth0 & 1.0.0.5/31 & r1-eth0.isp102.lab \\ 
r1 & eth1 & 1.102.3.0/31 & r1-eth1.isp102.lab \\ 
r1 & eth2 & 1.102.3.2/31 & r1-eth2.isp102.lab \\ 
r1 & eth3 & 1.102.3.5/31 & r1-eth3.isp102.lab \\ 
r1 & dummy0 & 1.102.4.1/32 & r1-lo.isp102.lab \\ 
\hline
r2 & eth0 & 2.21.0.1/31 & r2-eth0.isp102.lab \\ 
r2 & eth1 & 1.102.3.1/31 & r2-eth1.isp102.lab \\ 
r2 & eth2 & 1.102.3.6/31 & r2-eth2.isp102.lab \\ 
r2 & dummy0 & 1.102.4.2/32 & r2-lo.isp102.lab \\ 
\hline
r3 & eth0 & 1.102.1.1/24 & r3-eth0.isp102.lab \\ 
r3 & eth1 & 1.102.3.3/31 & r3-eth1.isp102.lab \\ 
r3 & eth2 & 1.102.3.7/31 & r3-eth2.isp102.lab \\ 
r3 & eth3 & 1.102.3.8/31 & r3-eth3.isp102.lab \\ 
r3 & dummy0 & 1.102.4.3/32 & r3-lo.isp102.lab \\ 
\hline
r4 & eth0 & 1.102.2.1/24 & r4-eth0.isp102.lab \\ 
r4 & eth1 & 1.102.3.4/31 & r4-eth1.isp102.lab \\ 
r4 & eth2 & 1.102.3.9/31 & r4-eth2.isp102.lab \\ 
r4 & dummy0 & 1.102.4.4/32 & r4-lo.isp102.lab \\ 
\hline
s1 & eth0 & 1.102.1.2/24 & ns.isp102.lab \\ 
s2 & eth0 & 1.102.1.3/24 & www.isp102.lab \\ 
s3 & eth0 & 1.102.1.4/24 & dhcpd.isp102.lab \\ 
\hline
c1 & eth0 & 1.102.2.X/24 ($X \in \{10,11\}$) & c1.isp102.lab \\ 
c2 & eth0 & 1.102.2.Y/24 ($Y \in \{10,11\}$) & c2.isp102.lab \\ 
\hline
\end{tabular}
\end{center}

\section{Routing and service implementation}
This section describes ISP implementation to realize routing and service requirements.

\subsection{Routing}
This section describes ISP implementation to fulfill routing requirements.

\subsubsection{Intra-domain routing}
We choose OSPF as our Interior Gateway Protocol (IGP) because of its fast convergence and flexible link metric configuration, which allows deterministic path selection without ECMP.

Our internal network connects four routers using five /31 point-to-point links. The router topology provides two edge-disjoint paths between each pair of routers, preserving internal connectivity after any single inter-router link failure once OSPF converges. To eliminate ECMP, both \texttt{r1}--\texttt{r4} and \texttt{r2}--\texttt{r3} are assigned an OSPF cost of 10, the core link \texttt{r1}--\texttt{r2} is assigned a cost of 5, while \texttt{r1}--\texttt{r3} has a cost of 9, and \texttt{r3}--\texttt{r4} a cost of 11.
The primary and secondary paths for the required traffic flows are listed below. Costs are sums of inter-router link costs; they exclude the common final LAN interface cost. Secondary paths assume one inter-router link on the primary path fails, with the remaining links operational:
\begin{itemize}
    \item \textbf{\texttt{r1} $\leftrightarrow$ clients (\texttt{r4})}: Primary is direct \texttt{r1}--\texttt{r4} (cost 10) for the lowest configured OSPF cost. Secondary is via \texttt{r3} (cost 20) upon link failure.
    
    \item \textbf{\texttt{r1} $\leftrightarrow$ servers (\texttt{r3})}: Primary is direct \texttt{r1}--\texttt{r3} (cost 9) for direct server access. Secondary is via \texttt{r2} (cost 15).
    
    \item \textbf{\texttt{r2} $\leftrightarrow$ clients (\texttt{r4})}: Primary is via \texttt{r1} (cost 15). Secondary is via \texttt{r3} (cost 21). There is no direct link between \texttt{r2} and \texttt{r4} in this topology.
    
    \item \textbf{\texttt{r2} $\leftrightarrow$ servers (\texttt{r3})}: Primary is direct \texttt{r2}--\texttt{r3} (cost 10). Secondary is via \texttt{r1} (cost 14) upon link failure.
    
    \item \textbf{clients (\texttt{r4}) $\leftrightarrow$ servers (\texttt{r3})}: Primary is direct \texttt{r4}--\texttt{r3} (cost 11) to keep local traffic off border routers. Secondary is via \texttt{r1} (cost 19).
\end{itemize}

Table~\ref{tab:primary_path} and Table~\ref{tab:secondary_path} below show the primary and secondary routing paths respectively.


% The table below shows the primary routing path between two devices, presented in a similar way as a bus fare matrix: each row represents a source with a path to each destination specified in each column, with each cell listing the intermediate path from the source to the destination.
\begin{table}[h!]
    \centering
    \begin{tabular}{|m{1.5cm}||m{2cm}|m{2cm}|m{2cm}|m{2cm}|}
        \hline
        \textbf{Path} & \textbf{r1} & \textbf{r2} & \textbf{servers} & \textbf{clients}\\ 
        \hline\hline
        \textbf{r1}         & X           & -           & r3             & r4\\
        \textbf{r2}         & -           & X           &r3               & r1 r4 \\ 
        \textbf{servers}    & r3           & r3           & X               & r3 r4 \\ 
        \textbf{clients}    & r4           & r4 r1          &r4 r3               & X \\ 
        \hline
    \end{tabular}
    \caption{Intermediate nodes in the primary routing path from row to column. X represents a path to itself, - represents a direct link without any intermediate node.}
    \label{tab:primary_path}
\end{table}

% The table below shows the secondary routing path between two devices, assuming the primary path fails.
\begin{table}[h!]
    \centering
    \begin{tabular}{|m{1.5cm}||m{2cm}|m{2cm}|m{2cm}|m{2cm}|}
        \hline
        \textbf{Path} & \textbf{r1} & \textbf{r2} & \textbf{servers} & \textbf{clients}\\ 
        \hline\hline
        \textbf{r1}         & X           & r3          & r2 r3              & r3 r4 \\
        \textbf{r2}         & r3          & X           & r1 r3              & r3 r4 \\ 
        \textbf{servers}    & r3 r2          & r3 r1          & X               & r3 r1 r4 \\ 
        \textbf{clients}    & r4 r3          & r4 r3          & r4 r1 r3              & X \\ 
        \hline
    \end{tabular}
    \caption{Intermediate nodes in the secondary routing path from row to column (when the primary routing path fails). 'X' represents a path to itself, '-' represents a direct link without any intermediate node.}
    \label{tab:secondary_path}
\end{table}



\subsubsection{Inter-domain routing}

We deploy BGP on border routers \texttt{r1} (peering with AS1) and \texttt{r2} (peering with AS21), connected via an iBGP session over loopback interfaces (\texttt{dummy0}) with \texttt{next-hop-self}. Core routers \texttt{r3} and \texttt{r4} run solely OSPF Area~0 without BGP. The configuration uses no \texttt{weight}, static routes, or IGP on external links. All routing policies rely on standard attributes (\texttt{LOCAL\_PREF}, \texttt{AS-PATH}, \texttt{COMMUNITY}), where \texttt{LOCAL\_PREF} is applied exclusively to incoming routes.

\paragraph{Prefix Aggregation and Outbound Traffic}

Both border routers originate their own dummy0 /32 into BGP and use \texttt{aggregate-address 1.102.0.0/20 summary-only} to advertise our aggregate while suppressing its more-specific BGP routes. Export route-maps control which additional prefixes may be announced: \texttt{TO-AS1} permits only our aggregate and AS21's \texttt{2.21.0.0/20}; \texttt{TO-AS21} permits our aggregate and the remaining eligible BGP routes with different prepend policies. Thus AS21 can use AS102 as backup transit, while AS102 does not export arbitrary AS21-learned prefixes to AS1. For outbound traffic, \texttt{r2} sets $LP = 200$ on incoming \texttt{2.21.0.0/20} from AS21 and $LP = 50$ on all other AS21 routes, while \texttt{r1} sets $LP = 100$ on routes from AS1. This directs AS21-bound traffic to exit via \texttt{r2} (higher LP) and regular Internet traffic via \texttt{r1}. If the \texttt{r1}--AS1 link fails, \texttt{r1} retains only iBGP-learned routes from \texttt{r2} (normally $LP = 50$ for other external destinations, and $LP = 200$ for AS21 itself), so outbound Internet traffic shifts to the backup path via AS21.

\paragraph{Inbound Traffic and Mutual Transit}

For incoming traffic, \texttt{r1} tags its \texttt{1.102.0.0/20} advertisement to AS1 with community \texttt{1:200}, prompting AS1 to set its internal LP to 200 and prefer the direct link; this community is applied exclusively to our own aggregated prefix and not to any transit routes. On the \texttt{r2}--AS21 peering, \texttt{r2} applies differentiated AS-path prepending: it prepends our ASN once for \texttt{1.102.0.0/20} and two times for all other (transit) routes. The single prepend on our own prefix makes AS21's direct path (length~2) shorter than the indirect AS2$\to$AS1 path (length~3), so AS21 prefers the direct link for reaching AS102; at the same time, AS2 sees the path via AS21 (length~3) as longer than via AS1 (length~2) and thus routes inbound traffic through AS1. The double prepend on transit routes ensures that AS21 strongly prefers its own AS2 link for Internet traffic during normal operation, and also prefers the AS2--AS3--AS1 backup path (length~3) over AS102 for AS1 destinations if the AS1--AS2 link fails. If the AS21--AS2 link fails, the AS102 route remains available as backup transit. Conversely, \texttt{r1} re-advertises AS21's prefix (\texttt{2.21.0.0/20}) to AS1 with two additional copies of ASN~102. Including the ASN automatically added on eBGP export, AS1 sees \texttt{102 102 102 21} (length~4). This is longer than both \texttt{2 21} (length~2) and the backup backbone path \texttt{3 2 21} (length~3), avoiding an equal-length tie when AS1--AS2 fails. If AS21--AS2 fails, AS1 can still reach AS21 through AS102.


\paragraph{Internal Routing Consistency and Loop Prevention}
Core routers \texttt{r3} and \texttt{r4} use OSPF instead of BGP. Border routers originate conditional OSPF defaults with external Type-1 metrics of 10 (\texttt{r1}) and 50 (\texttt{r2}), preferring \texttt{r1} for general external traffic. Each \texttt{default-information originate always} command has a route-map: \texttt{DIRECT-AS1} requires \texttt{1.0.0.0/20} with next hop \texttt{1.0.0.4}, while \texttt{ONLY-AS21} requires \texttt{2.21.0.0/20} with next hop \texttt{2.21.0.0}. These maps also filter BGP redistribution into OSPF. The \texttt{always} keyword removes the need for an installed default route, but does not remove the route-map condition. An alternative iBGP route does not satisfy the direct-peer next-hop match.

When the direct AS1 route is withdrawn after an uplink or eBGP-session failure, \texttt{r1} withdraws its default LSA and its redistributed AS1 prefix. After convergence, the core routers use \texttt{r2}'s default. If \texttt{r1}--\texttt{r2} also fails, \texttt{r1}'s remaining iBGP route can forward via \texttt{r3} to \texttt{r2}, while \texttt{r3} no longer defaults back to \texttt{r1}. If neither direct upstream route remains, neither border router originates a default.

For AS21 traffic, the more-specific redistributed \texttt{2.21.0.0/20} sends core traffic towards \texttt{r2}, including when \texttt{r1}--\texttt{r2} fails. Both border routers also apply \texttt{ip protocol ospf route-map OSPF-TO-KERNEL}: it denies only the exact OSPF default route from being installed in the kernel forwarding table, and permits other OSPF routes. This prevents a border router from following the other border's OSPF default back through a core router that forwards towards it; default LSAs remain available to the core. These mechanisms address the observed forwarding loops after convergence; failure detection and route withdrawal can still cause temporary packet loss.

\subsection{Internet service}
This section describes ISP implementation to fulfil service requirements.


\subsubsection{DNS}

% \todo[inline]{Describe which device you use as the DNS server and what IP address and name you assign to it.
% You need to ensure that all clients and servers can resolve the names of all servers on the Internet. Similarly, you need to make sure that other ASes can resolve the names of devices in your ISP. Briefly explain what you do to ensure your DNS service works as expected.}

The DNS server, represented as 's1' in Figure \ref{fig:your_figure} has the IP address \textbf{1.102.1.2}, and the name \textbf{ns.isp102.lab}.

BIND~9 on \texttt{s1} provides both recursive resolution for AS102 and authoritative service for \texttt{isp102.lab} and \texttt{102.1.in-addr.arpa}. Clients and servers use \texttt{1.102.1.2} as their DNS resolver. Recursion is allowed for \texttt{1.102.0.0/20}; authoritative queries to our two zones are allowed from any address.

For external names, the resolver starts from its root hints at \texttt{1.0.1.2} and follows DNS referrals. The lab's \texttt{.lab} zone delegates \texttt{isp102.lab} to \texttt{ns.isp102.lab}, with address \texttt{1.102.1.2}. For reverse resolution, the lab hierarchy delegates \texttt{1.in-addr.arpa} to AS1's DNS server, which delegates \texttt{102.1.in-addr.arpa} to our server. Our reverse zone contains PTR records for internal router interfaces, loopbacks, servers, and the two configured client addresses. The external-facing addresses \texttt{1.0.0.5} and \texttt{2.21.0.1} are outside this reverse zone.

The static client A/PTR records map \texttt{c1} to \texttt{1.102.2.10} and \texttt{c2} to \texttt{1.102.2.11}. DHCP uses a dynamic pool without fixed reservations or DNS updates, so these names are not guaranteed to identify the corresponding client after leases are reassigned.

\subsubsection{Web}
% \todo[inline]{Describe which device you use as the web server and what IP address and name you assign to it.}

The Web server is denoted as 's2' in Figure \ref{fig:your_figure} and has IP address \textbf{1.102.1.3}, with name \textbf{www.isp102.lab}. The web server also has a main page served from \texttt{index.html} with the following contents:

\begin{itemize}
    \item  ASN: 102
\item NETWORK: 1.102.0.0/20
\item NAME1: Emanuel Paraschiv
\item EMAIL1: emapar@kth.se
\item NAME2: Haoyu Chen 
\item EMAIL2: chchen3@kth.se 
\end{itemize}


\subsubsection{DHCP}
% \todo[inline]{Describe which device you use as the DHCP server and DHCP relay and what IP addresses and names they have. Briefly explain how you set up the DHCP server and DHCP relay.}

The DHCP server is represented as 's3' in Figure \ref{fig:your_figure}, having the IP address \textbf{1.102.1.4} and name \textbf{dhcpd.isp102.lab}.


Given that clients 'c1' and 'c2' live on a different subnet
(\textbf{1.102.2.0/24}), the gateway \textbf{r4} runs isc-dhcp-relay and unicasts requests to \textbf{1.102.1.4}, added to its existing routing functionalities. 


For the client LAN (\textbf{1.102.2.0/24}), the DHCP server hands out addresses dynamically from \textbf{1.102.2.10}--\textbf{1.102.2.11} (two addresses). It supplies gateway \texttt{1.102.2.1}, DNS server \texttt{1.102.1.2}, domain \texttt{isp102.lab}, and subnet mask \texttt{255.255.255.0}. The default and maximum lease times are 1000 and 5000 seconds, respectively.




\end{document}
